\documentclass[11pt]{article}
\usepackage[a4paper,margin=26mm]{geometry}
\usepackage{amsmath,amssymb,amsthm,booktabs,array,longtable}
\usepackage[T1]{fontenc}
\usepackage[colorlinks=true,linkcolor=blue,citecolor=blue,urlcolor=blue]{hyperref}
\newtheorem{theorem}{Theorem}
\newtheorem{lemma}[theorem]{Lemma}
\newtheorem{corollary}[theorem]{Corollary}
\newcommand{\FF}{\mathbb F_2}
\newcommand{\ZZ}{\mathbb Z}
\newcommand{\RR}{\mathbb R}
\newcommand{\ep}{\varepsilon}
\newcommand{\leanbase}{https://github.com/the-omega-institute/trureturing/blob/42aeb4e95344fb1ea9d2042d02ff1a5589faaacb/D5/S3/Combinatorics/Hankel/}
\title{Inertia extensions of Krattenthaler's identity for Sun's Hankel determinant conjecture}
\author{Haobo Ma\thanks{ChronoAI Pte Ltd, Singapore.}\and
Wenlin Zhang\thanks{National University of Singapore; ChronoAI Pte Ltd, Singapore.}}
\date{November 13, 2026}
\begin{document}
\maketitle
\begin{abstract}
For Sun's \emph{On some determinants with Legendre symbol entries},
arXiv:1308.2900v5, Conjecture 3.16, p.~17, we extend the zero-class
identity credited to C.\ Krattenthaler in arXiv:1308.2900v6
(April 24, 2014), Remark 3.16.
For the sequence A006077, the Hankel determinant of order $n+1$ vanishes
exactly when $n\equiv1\pmod3$; otherwise its quotient by
$(-1)^{\lfloor(n+1)/3\rfloor}6^n$ is a positive odd integer.
Binomial translation gives a three-block factorization. Integral row
normalization reduces parity to the Rueppel determinants, and a positive
moment representation gives the sign. We also give parameter-independent
and square-class consequences, the exact inertia, and a two-parameter
deformation. Related moment arguments establish the harmonic sign assertion
in Conjecture 3.9(i) and the two squared-number assertions in Conjecture 3.17.
\end{abstract}

\section{Statement and attribution}
For nonnegative integers $r$ and $n$, put
\[
 c_r=\binom{2r}{r}\binom{3r}{r}=\frac{(3r)!}{(r!)^3},\qquad
 w_n=\sum_{r=0}^{\lfloor n/3\rfloor}(-1)^r3^{n-3r}\binom n{3r}c_r,
 \qquad W_n=\det(w_{i+j})_{0\leq i,j\leq n}.
\]
An empty determinant is $1$. Sun posed the following assertion on August 21,
2013~\cite[Conjecture 3.16, p.~17]{sun5}; the same assertion is recorded in
OEIS A006077~\cite{oeis}.
\begin{theorem}\label{thm:main}
For every $n\geq0$,
\[
 W_n=0\quad\Longleftrightarrow\quad n\equiv1\pmod3.
\]
If $n\not\equiv1\pmod3$, then
\[
 \frac{(-1)^{\lfloor(n+1)/3\rfloor}W_n}{6^n}
 \quad\hbox{is a positive odd integer.}
\]
\end{theorem}
Sun's next version credits the identity $W_{3m+1}=0$ to
C.\ Krattenthaler~\cite[Remark 3.16]{sun6}. Thus that identity is an earlier
result, not a new consequence attributed here to the present authors.
The nonzero-class assertion also appears in Sun's 2019 survey,
Conjecture 81(i), p.~68, with the zero class attributed separately in
Remark 81\cite{sun2019}.
Our argument combines this residue-block structure with the positivity and
parity of two multinomial Hankel families. The binary determinant ingredients
are given by Bacher's triangular factorizations\cite[Theorems 1.2(i) and 10.1(i)]{bacher}
and later by Cigler~\cite[Theorems 1.1 and 2.1]{cigler}. We give the needed reductions
explicitly. The weighted binomial invariance used below is the falling
$k$-binomial invariance of Spivey and Steil~\cite[Theorem 6.1]{spivey}.
The positive moment representation of $c_r$ is given by Penson and
Solomon~\cite[Example 9]{penson}; Section 4 derives the needed positivity
from a product of beta variables.
Later permutation proofs of the ordinary and shifted binary ingredients
are given, respectively, by Grinberg (January 9, 2018) and
Random (July 19, 2026)\cite{community}.
For comparison, Zhu and Sun~\cite[Theorems 1.1 and 1.2]{zhusun} prove
positive odd normalized Hankel determinants for Franel, Domb, and
Catalan--Larcombe--French numbers. Those results settle other conjectures
from the same source, but concern different sequences from A006077.
Their general-Franel theorem proves odd integrality, not positivity for
every order; their Ap\'ery theorem likewise proves integrality alone.
Edgar's positive moment density for the sequence $A_n$\cite[Theorem 1]{edgar}
supplies its positivity assertion, while Adamczewski\cite[Table 1, p.~16,
A228304]{adamczewski} gives the two congruences of Conjecture 3.12.

\section{Translation and the three residue classes}
Define a sparse sequence by
\[
 b_{3r}=(-1)^r c_r,\qquad b_{3r+1}=b_{3r+2}=0.
\]
Then $w_n=\sum_{j=0}^n\binom nj3^{n-j}b_j$.
More generally, if $\widetilde b_n=\sum_{j=0}^n\binom nj a^{n-j}b_j$,
the lower triangular matrix
$P_{ij}=\binom ij a^{i-j}$ for $j\leq i$ (and $0$ otherwise) satisfies
\[
 (\widetilde b_{i+j})_{0\leq i,j<N}
   =P(b_{i+j})_{0\leq i,j<N}P^{\mathsf T}.
\]
Indeed, the coefficient of $a^{i+j-t}b_t$ on the right is
$\sum_{u+v=t}\binom iu\binom jv=\binom{i+j}t$, by multiplying
$(1+x)^i(1+x)^j$. Since $P$ has diagonal $1$,
\begin{equation}\label{eq:translate}
 W_n=\det(b_{i+j})_{0\leq i,j\leq n}.
\end{equation}
For $k\geq0$ set
\[
 D_k=\det(c_{i+j})_{0\leq i,j<k},\qquad
 E_k=\det(c_{i+j+1})_{0\leq i,j<k}.
\]
\begin{lemma}\label{lem:blocks}
For $m\geq0$,
\begin{align}
 W_{3m}&=(-1)^mD_{m+1}E_m^2,\label{eq:block0}\\
 W_{3m+1}&=0,\label{eq:block1}\\
 W_{3m+2}&=(-1)^{m+1}D_{m+1}E_{m+1}^2.\label{eq:block2}
\end{align}
\end{lemma}
\begin{proof}
Multiply row $i$ and column $i$ of the sparse matrix by
$(-1)^{\lfloor i/3\rfloor}$. The determinant is unchanged because the two
diagonal multipliers have product $1$. Group the indices by residues
$0,1,2$ modulo $3$, using the same permutation for rows and columns.
For indices $3p+r,3q+s$, the nonzero blocks have entries $c_{p+q}$ if
$r=s=0$, and $-c_{p+q+1}$ if $r+s=3$. All other blocks are zero.

For $n=3m$, the residue sizes are $(m+1,m,m)$, so the resulting matrix is
the direct sum of $(c_{p+q})_{p,q<m+1}$ and
$\left(\begin{smallmatrix}0&-C\\-C&0\end{smallmatrix}\right)$,
where $C=(c_{p+q+1})_{p,q<m}$. Swapping the two column blocks gives
determinant $(-1)^m\det(C)^2$, proving \eqref{eq:block0}.
For $n=3m+2$ the three sizes are $(m+1,m+1,m+1)$ and the same calculation
proves \eqref{eq:block2}.
For $n=3m+1$ the sizes are $(m+1,m+1,m)$. Every row in the residue-$1$
class is supported on the $m$ columns in the residue-$2$ class. These
$m+1$ rows are linearly dependent, proving \eqref{eq:block1}.
\end{proof}

\section{Integral normalization and parity}
For $r\geq1$, define the integer
\[
 q_r=\binom{2r-1}{r-1}\binom{3r-1}{r-1}.
\]
The two binomial identities
$\binom{2r}r=2\binom{2r-1}{r-1}$ and
$\binom{3r}r=3\binom{3r-1}{r-1}$ give $c_r=6q_r$.
Define
\[
 d_m=\det(A_{ij})_{0\leq i,j\leq m},\quad
 A_{0j}=c_j,\quad A_{ij}=q_{i+j}\ (i\geq1),\qquad
 e_m=\det(q_{i+j+1})_{0\leq i,j<m}.
\]
Dividing all but the first row in $D_{m+1}$, and every row in $E_m$, by
$6$ proves
\begin{equation}\label{eq:normalize}
 D_{m+1}=6^m d_m,\qquad E_m=6^m e_m.
\end{equation}

\begin{lemma}\label{lem:parity}
Both $d_m$ and $e_m$ are odd for every $m\geq0$.
\end{lemma}
Zhu and Sun's Lemma 2.3\cite{zhusun}, applied with their $k=3$ and
$x_r=c_r$, already gives the oddness of $D_{m+1}/6^m$ once the
coefficient parity below is known. We include the binary argument for
both normalized families. Its recurrences are the characteristic-$2$
specializations of Guo's Lemma 2(3)--(6)\cite{guo}.
\begin{proof}
Let $s_2(r)$ count the ones in the binary expansion of $r$, and let
$\nu_2$ denote the exponent of $2$ in a nonzero integer.
The identity
$\nu_2(r!)=\sum_{j\geq1}\lfloor r/2^j\rfloor=r-s_2(r)$ gives
\[
 \nu_2\binom{2r}r=s_2(r),\qquad
 \nu_2\binom{3r}r=s_2(r)+s_2(2r)-s_2(3r).
\]
Consequently, $q_r$ odd implies $s_2(r)=1$, since $c_r=6q_r$ has
valuation $1$ and the two binomial valuations are nonnegative. Conversely,
if $r=2^j$, these valuations are $1$ and $0$, respectively, so $q_r$ is odd.

Work now in $\FF$, and let $\ep(r)$ be $1$ if $r$ is a positive power of
$2$, and $0$ otherwise, including $r=0$. Thus
\[
 q_r\equiv\ep(r)\pmod2\ (r\geq1),\qquad
 \ep(2r)=\ep(r),\qquad
 \ep(2r+1)=\begin{cases}1&r=0,\\0&r>0.\end{cases}
\]
Put
\[
 S_k=\det(\ep(i+j+1))_{0\leq i,j<k},\qquad
 T_k=\det(\ep(i+j+2))_{0\leq i,j<k}
 \quad\hbox{in }\FF.
\]
Separating even and odd indices gives
\begin{equation}\label{eq:rueppel}
 S_{2k}=S_k^2,\qquad S_{2k+1}=T_k^2,\qquad
 T_N=S_{\lceil N/2\rceil}T_{\lfloor N/2\rfloor}.
\end{equation}
Here are the determinant reductions, including the odd-size case.
The matrix for $S_{2k}$ becomes
$\left(\begin{smallmatrix}J&B\\B&0\end{smallmatrix}\right)$,
where $J$ has its only possible nonzero entry $J_{00}=1$, and
$B=(\ep(i+j+1))_{i,j<k}$. Swapping the two column blocks gives
determinant $\det(B)^2$ in characteristic $2$.
For $S_{2k+1}$ order the indices as $0$; $2,4,\ldots,2k$;
$1,3,\ldots,2k-1$. With
$C=(\ep(i+j+2))_{i,j<k}$ and $u_j=\ep(j+1)$ the matrix is
\[
 \begin{pmatrix}1&0&u\\0&0&C\\u^{\mathsf T}&C&0\end{pmatrix}.
\]
Swapping the last two column groups makes this block triangular with
diagonal blocks $\left(\begin{smallmatrix}1&u\\0&C\end{smallmatrix}\right)$
and $C$, proving the second identity.
Finally the matrix for $T_N$ becomes block diagonal with blocks for
$S_{\lceil N/2\rceil}$ and $T_{\lfloor N/2\rfloor}$, because all
off-diagonal entries are $\ep(2r+1)$ with $r\geq1$.

The initial values are $S_0=T_0=S_1=T_1=1$.
For $N\geq2$ every index on the right of \eqref{eq:rueppel} is smaller
than $N$. Strong induction therefore proves $S_N=T_N=1$ for all $N$.
The determinant $e_m$ reduces to $S_m$. The first row of $A$ reduces
to $(1,0,\ldots,0)$, since $c_0=1$ and $6\mid c_j$ for $j>0$.
Expansion along that row shows $d_m\equiv T_m=1\pmod2$ as well.
\end{proof}

\section{Positive moments and completion of the proof}
\begin{lemma}\label{lem:positive}
For all $k\geq0$, $D_k>0$ and $E_k>0$.
\end{lemma}
\begin{proof}
For $0<a<1$ let $U_a$ have probability density
\[
 \frac{t^{a-1}(1-t)^{-a}}{B(a,1-a)}\quad(0<t<1),\qquad
 B(a,1-a)=\int_0^1t^{a-1}(1-t)^{-a}\,dt.
\]
The integral is finite and positive. Integration by parts in the beta
integrals gives
$\mathbb E(U_a^r)=(a)_r/r!$, where
$(a)_r=a(a+1)\cdots(a+r-1)$ and $(a)_0=1$.
Explicitly, integration of the derivative of
$t^{a+r}(1-t)^{1-a}$, whose endpoint values are zero, gives
$B(a+r+1,1-a)=\frac{a+r}{r+1}B(a+r,1-a)$; induction gives the formula.
Take independent $U_{1/3}$ and $U_{2/3}$ and set
$X=27U_{1/3}U_{2/3}$. Grouping the factors of $(3r)!$ by residues modulo
$3$ proves
\[
 \mathbb E(X^r)=27^r\frac{(1/3)_r(2/3)_r}{(r!)^2}=c_r.
\]
Every open subinterval of $(0,27)$ has positive probability for $X$:
its inverse image under $(u,v)\mapsto27uv$ contains a nonempty open
rectangle in $(0,1)^2$, where both densities are positive.
For any nonzero real polynomial $p$ and any integer $s\geq0$,
$\mathbb E[X^s p(X)^2]>0$. Indeed, a nonzero polynomial has finitely
many zeros, so some such interval has $X^s p(X)^2$ bounded below by a
positive constant. Hence every matrix $(c_{i+j+s})_{i,j<k}$ is positive
definite for $k\geq1$, by applying this inequality to
$p(x)=\sum_{i<k}v_i x^i$ for $v\ne0$. Its determinant is positive.
This proves the assertion for $s=0,1$; the empty determinants equal $1$.
\end{proof}
By \eqref{eq:normalize}, $d_m,e_m>0$, and Lemma~\ref{lem:parity} makes
them positive odd integers. Substitution in Lemma~\ref{lem:blocks} yields
\begin{align}
 \frac{(-1)^mW_{3m}}{6^{3m}}&=d_m e_m^2,\label{eq:quot0}\\
 \frac{(-1)^{m+1}W_{3m+2}}{6^{3m+2}}&=d_m e_{m+1}^2.\label{eq:quot2}
\end{align}
These prove Theorem~\ref{thm:main}, including the nonvanishing and the
case $n=0$, where $d_0=e_0=W_0=1$.

\begin{corollary}\label{cor:extra}
For $n\not\equiv1\pmod3$, $\nu_2(W_n)=n$.
The positive normalized determinants at $n=3m$ and $n=3m+2$ represent the
same square class in $\mathbb Q_{>0}^{\times}/(\mathbb Q_{>0}^{\times})^2$,
namely the class of $d_m$.
Moreover, replacing $3^{n-3r}$ in the definition of $w_n$ by
$a^{n-3r}$ for any $a\in\RR$ leaves every $W_n$ unchanged.
\end{corollary}
\begin{proof}
The first two assertions follow directly from \eqref{eq:quot0} and
\eqref{eq:quot2}. The last follows from the triangular identity
\eqref{eq:translate}, which holds for every real $a$.
\end{proof}

\begin{corollary}\label{cor:inertia}
For real $a,t$ define
\[
 w_j(a,t)=\sum_{r=0}^{\lfloor j/3\rfloor}
  \binom j{3r}a^{j-3r}t^r c_r,\qquad
 H_n(a,t)=(w_{i+j}(a,t))_{0\leq i,j\leq n}.
\]
For $t\ne0$, the inertia of $H_n(a,t)$, listed as the numbers of positive,
negative and zero eigenvalues, is
\[
\begin{array}{c|c}
n&\operatorname{inertia} H_n(a,t)\\\hline
3m&(2m+1,m,0)\\
3m+1&(2m+1,m,1)\\
3m+2&(2m+2,m+1,0).
\end{array}
\]
In particular, every zero in the original determinant family has nullity
exactly one. For $n\not\equiv1\pmod3$,
\[
 \det H_n(a,t)=t^{n(n+1)/3}W_n.
\]
For $t=0$, the inertia is $(1,0,n)$ and the determinant is $1$ at $n=0$
and $0$ otherwise. For nonzero integers $t$ and $n\not\equiv1\pmod3$,
\[
 \nu_2(\det H_n(a,t))=n+\frac{n(n+1)}3\nu_2(t)
 \quad\text{when }a\in\ZZ.
\]
\end{corollary}
\begin{proof}
Translation is a real congruence by the invertible matrix $P$, so it
preserves inertia. The sparse matrix with entries $c_{(i+j)/3}$ when
$3\mid i+j$, and zero otherwise, has the positive definite block
from $D$ and the paired block
$\left(\begin{smallmatrix}0&C\\C^{\mathsf T}&0\end{smallmatrix}\right)$.
In the two balanced cases $C$ is positive definite of size $m$ or $m+1$;
the change of variables $(x,y)\mapsto(x+y,x-y)$ gives one positive and
one negative direction for each dimension. In the unbalanced case $C$
is $(m+1)\times m$ and has full column rank because its first $m$ rows
form the positive definite matrix for $E_m$. Invertible row and column
changes reduce it to $\left(\begin{smallmatrix}I_m\\0\end{smallmatrix}\right)$,
so the paired block has inertia $(m,m,1)$. This proves the table, also
when $m=0$.

Let $\rho$ be the real cube root of $t\ne0$. Multiplication on both sides
by $\operatorname{diag}(1,\rho,\ldots,\rho^n)$ changes the sparse
$t=1$ matrix into the sparse $t$ matrix. Its determinant multiplier is
$\rho^{n(n+1)}=t^{n(n+1)/3}$ in the nonzero classes. This exponent is
even: it is $m(3m+1)$ or $(m+1)(3m+2)$. Taking $t=-1$ identifies the
base determinant with $W_n$. The valuation follows from
Corollary~\ref{cor:extra}. At $t=0$ the entries are $a^{i+j}$, the
outer product of $(1,a,\ldots,a^n)$ with itself.
\end{proof}

\section{Moment consequences for two sibling conjectures}
The assertions below concern different sequences, not additional parts of
the A006077 assertion. They use the same positivity principle: a measure
with interval support makes polynomial Gram matrices strictly positive.

\begin{corollary}\label{cor:harmonic}
For every real $\alpha>0$ and integer $n\geq1$, with
$h_j=\sum_{k=1}^j k^{-\alpha}$ and $h_0=0$,
\[
 (-1)^n\det(h_{i+j})_{0\leq i,j\leq n}>0.
\]
Thus Sun's v5 Conjecture 3.9(i), p.~15, holds for every positive integer
order, and also for every positive real order. For order one, a prior
exact evaluation follows from Petrov's answer of February 16, 2017,
Cigler's completion of February 17, and Amdeberhan's independent identity
proof of February 18; see also Cigler, Theorem 4.1, equation (4.2)\cite{harmonic}.
\end{corollary}
\begin{proof}
The substitution $x=e^{-u}$ in the gamma integral gives
\[
 h_j=\int_0^1\frac{1-x^j}{1-x}\,
       \frac{(-\log x)^{\alpha-1}}{\Gamma(\alpha)}\,dx.
\]
Subtract row $0$ from each other row and column $0$ from each other
column. The resulting matrix is
$\left(\begin{smallmatrix}0&h^{\mathsf T}\\h&-G\end{smallmatrix}\right)$,
where $h=(h_1,\ldots,h_n)^{\mathsf T}$ and
\[
 G_{ij}=\int_0^1 (1-x^i)(1-x^j)
 \frac{(-\log x)^{\alpha-1}}{\Gamma(\alpha)(1-x)}\,dx
 \quad(1\leq i,j\leq n).
\]
The integrals converge at both endpoints. The polynomials $1-x^i$
are linearly independent and the weight is positive on $(0,1)$; hence
$G$ is positive definite. The Schur complement of $-G$ is
$h^{\mathsf T}G^{-1}h>0$ since $h_1=1$. Therefore the inertia is
$(1,n,0)$, proving the sign.
\end{proof}

\begin{corollary}\label{cor:squares}
Let $B_j$ and $E_j$ denote the Bernoulli and Euler numbers in the
conventions
\[
 \frac{z}{e^z-1}=\sum_{j\geq0}B_j\frac{z^j}{j!},\qquad
 \frac1{\cosh z}=\sum_{j\geq0}E_j\frac{z^j}{j!}.
\]
For every integer $n\geq1$,
\[
 \det(B_{i+j}^2)_{0\leq i,j\leq n}<0,
 \qquad \det(E_{i+j}^2)_{0\leq i,j\leq n}>0.
\]
These are the two assertions of Sun's v5 Conjecture 3.17, p.~18.
\end{corollary}
\begin{proof}
Let $U,V$ be independent exponential variables of mean one, and let
$\epsilon$ be an independent uniform sign. Write $\mu$ for the positive
measure $4\sum_{k,l\geq1}\mathcal L(\epsilon UV/(4\pi^2kl))$.
Its absolute moments of every integer order $j\geq2$ are finite.
For even $j\geq2$, Euler's zeta formula\cite[25.6.2]{dlmf} gives
\[
 \int x^j\,d\mu(x)
 =\frac{4(j!)^2\zeta(j)^2}{(2\pi)^{2j}}=B_j^2;
\]
the odd moments of orders at least three are zero. Although $\mu$ has
infinite total mass, the measure $x^2d\mu(x)$ is finite, has moments
of every order, and gives positive mass to every open interval.
Consequently $C=(B_{i+j}^2)_{1\leq i,j\leq n}$ is positive definite.
Write $b=(B_1^2,\ldots,B_n^2)^{\mathsf T}$. Its bordered matrix has
Schur complement $1-b^{\mathsf T}C^{-1}b$. The variational identity
\[
 b^{\mathsf T}C^{-1}b=\max_v(2b^{\mathsf T}v-v^{\mathsf T}Cv)
 \geq 2\cdot\tfrac14\cdot9-\tfrac1{36}\cdot9^2=\tfrac94
\]
uses $v=(9,0,\ldots,0)^{\mathsf T}$. Thus the complement is negative
and the Bernoulli-square determinant is negative.

For Euler numbers, let $Y$ have density $1/(2\cosh(\pi y/2))$ on
$\RR$. With $u=e^{\pi y}$, Euler's beta integral gives, for
$|z|<\pi/2$,
\[
 \mathbb E e^{zY}
 =\frac1\pi\int_0^\infty\frac{u^{z/\pi-1/2}}{1+u}\,du
 =\sec z.
\]
These exponential moments justify differentiating at zero. If $Z$ is
an independent copy of $Y$, the Taylor expansion of $\sec z$
\cite[4.19.5]{dlmf} gives
$\mathbb E(YZ)^j=(\mathbb E Y^j)^2=E_j^2$, including odd $j$.
The product $YZ$ has positive mass in every open interval of $\RR$.
Its polynomial Gram matrices are therefore positive definite, proving
the Euler-square inequality.
\end{proof}

\section{Coverage and verification}
All numbered items below refer to Sun v5 unless another version is given.
The labels apply to the precise assertions at those locators; the
inertia and deformation statements are not posed in the source.
\small
\begin{longtable}{@{}p{0.39\linewidth}p{0.56\linewidth}@{}}
\toprule
Source item & Coverage \\
\midrule
3.1--3.8 (each), pp.~12--15 & Not covered: character-symbol and reciprocal matrices, not these moment sequences. \\
3.9(i), p.~15 & Proved: Corollary~\ref{cor:harmonic}; order one reproves the Petrov--Cigler evaluation (2017)\cite{harmonic}. Positive real orders are not posed; proved here. \\
3.9(ii), p.~15 & Not covered: the harmonic determinant's prime congruence. \\
3.10(i), p.~15, ordinary Franel part & Prior proof cited: Zhu--Sun, Theorem 1.1\cite{zhusun}. \\
3.10(i), generalized Franel part & Proved for $r=2$ below. Prior odd integrality for $r\geq3$ cited, not general positivity; positivity for $r\geq4$ is not covered. \\
3.10(ii), p.~15 & Not covered: prime congruence. \\
3.11(i), p.~16 & Prior integrality cited: Zhu--Sun, Theorem 1.3. The $A_n$ part is positive by Edgar's Theorem 1; positivity of $b_n$ is not covered. \\
3.11(ii), p.~16 & Not covered: prime congruence. \\
3.12, p.~16 & Prior proof cited: Adamczewski (2026), Table 1, p.~16, A228304\cite{adamczewski}; both prime congruences. \\
3.13(i), p.~16; 3.14(i), p.~17 & Prior proofs cited: Zhu--Sun, Theorem 1.2. \\
3.13(ii), p.~16; 3.14(ii), p.~17 & Not covered: prime-square congruences. \\
3.15(i),(ii), p.~17 & Not covered: Catalan-binomial determinants. \\
3.16, p.~17 & Proved in full: Theorem~\ref{thm:main}; zero class reproved with prior attribution to Krattenthaler, attested by Sun, April 24, 2014. \\
3.17 (both inequalities), p.~18 & Proved: Corollary~\ref{cor:squares}. \\
v6, Remark 3.16, p.~18 & Reproved; prior result Krattenthaler, attested 2014. \\
OEIS A006077, August 21, 2013 comment & Proved: the same all-index assertion. \\
v6, Conjecture 3.16, p.~18; OEIS August 23, 2013 comment & Not covered: the later prime-square congruence. \\
Inertia, nullity, scale deformation, valuation, square classes & Not posed; proved in Corollaries~\ref{cor:extra} and~\ref{cor:inertia}. \\
\bottomrule
\end{longtable}
\normalsize
For completeness, the order-$2$ part of 3.10(i) follows by the same
normalization and moment argument. Vandermonde's identity gives
$f_j^{(2)}=\binom{2j}j=\mathbb E(4U)^j$ for
$U\sim\operatorname{Beta}(1/2,1/2)$. For $j>0$, write
$\binom{2j}j=2q'_j$ with $q'_j=\binom{2j-1}{j-1}$.
Its $2$-adic valuation is $s_2(j)-1$, so $q'_j$ is odd exactly at
powers of two. Dividing all but the first row of the determinant by
$2$ gives an integer determinant reducing modulo $2$ to $T_n=1$.
The interval support of $4U$ gives strict positivity. Hence
$2^{-n}\det(f_{i+j}^{(2)})_{0\leq i,j\leq n}$ is a positive odd integer.
For the $A_n$ part of 3.11(i), Edgar's strictly positive density on
$(0,17+12\sqrt2)$ makes every polynomial Gram matrix positive definite.
Combining this with Zhu--Sun's integral $24^{-n}$ normalization proves
the positive-integer assertion for that sequence.

The complete written proof is Sections 2--4. The following pinned Lean
files give the corresponding declarations, with namespace
\texttt{D5.S3.Combinatorics.Hankel.SunA006077}:
\begin{itemize}
\raggedright
\item \href{\leanbase SunA006077.lean}{\texttt{SunA006077.lean}}:
\texttt{source\_as\_translate} and \texttt{translate\_hankel\_det}
identify the literal sum and its Hankel invariance.
\item \href{\leanbase SunA006077Blocks.lean}{\texttt{SunA006077Blocks.lean}}:
\texttt{block\_zero}, \texttt{zero\_branch}, and \texttt{block\_two}
give all three residue classes.
\item \href{\leanbase SunA006077Binary.lean}{\texttt{SunA006077Binary.lean}}:
\texttt{rueppel\_determinants} proves both binary families equal $1$.
\item \href{\leanbase SunA006077Arithmetic.lean}{\texttt{SunA006077Arithmetic.lean}}:
\texttt{c\_six}, \texttt{q\_odd\_iff}, \texttt{D\_normalization},
\texttt{E\_normalization}, and \texttt{normalized\_odd} give integral
normalization and oddness.
\item \href{\leanbase SunA006077Positivity.lean}{\texttt{SunA006077Positivity.lean}}:
\texttt{multinomial\_posSemidef} gives positive semidefiniteness for every
nonnegative shift. The strict-support argument above gives positive
definiteness directly.
\item \href{\leanbase SunA006077Result.lean}{\texttt{SunA006077Result.lean}}:
\texttt{result} proves the complete all-index statement, obtaining strict
positivity from semidefiniteness and odd, hence nonzero, determinants.
\end{itemize}
Corollaries~\ref{cor:extra}--\ref{cor:squares}, including the real
parameters, inertia and sibling inequalities, are written proofs;
the pinned declarations concern the original A006077 family.
The accompanying
\texttt{verify.py} checks exact finite instances independently from the
literal sum; these finite checks complement, rather than replace, the
all-index proof.

\small
\begin{thebibliography}{99}
\bibitem{sun5} Z.-W. Sun,
\emph{On some determinants with Legendre symbol entries},
\href{https://arxiv.org/abs/1308.2900v5}{arXiv:1308.2900v5}, August 22, 2013,
Conjecture 3.16, p.~17.
\bibitem{sun6} Z.-W. Sun,
\emph{On some determinants with Legendre symbol entries},
\href{https://arxiv.org/abs/1308.2900v6}{arXiv:1308.2900v6},
April 24, 2014, Conjecture and Remark 3.16.
\bibitem{oeis} OEIS Foundation Inc.,
\href{https://oeis.org/A006077}{The On-Line Encyclopedia of Integer Sequences,
A006077}, comments and formula by Z.-W. Sun dated August 21 and 23, 2013,
accessed November 13, 2026.
\bibitem{cigler} J. Cigler,
\emph{A curious class of Hankel determinants},
\href{https://arxiv.org/abs/1803.05164v2}{arXiv:1803.05164v2} (2018),
Theorems 1.1 and 2.1.
\bibitem{spivey} M. Z. Spivey and L. L. Steil,
\emph{The $k$-binomial transforms and the Hankel transform},
Journal of Integer Sequences \textbf{9} (2006), Article 06.1.1, Theorem 6.1.
\href{https://cs.uwaterloo.ca/journals/JIS/VOL9/Spivey/spivey7.html}{Full text}.
\bibitem{penson} K. A. Penson and A. I. Solomon,
\emph{Coherent states from combinatorial sequences},
\href{https://arxiv.org/abs/quant-ph/0111151v1}{arXiv:quant-ph/0111151v1}
(2001), equation (6) and Example 9.
\bibitem{zhusun} B.-X. Zhu and Z.-W. Sun,
\emph{Hankel-type determinants for some combinatorial sequences},
International Journal of Number Theory \textbf{14} (2018), no.~5, 1265--1277.
\href{https://doi.org/10.1142/S1793042118500793}{doi:10.1142/S1793042118500793};
\href{https://arxiv.org/abs/1609.06810v4}{arXiv:1609.06810v4}.
\bibitem{harmonic} J. Cigler,
\emph{Hankel determinants of harmonic numbers and related topics},
\href{https://homepage.univie.ac.at/johann.cigler/preprints/harmonic.pdf}{author-hosted preprint},
Theorem 4.1, equation (4.2), p.~11;
F. Petrov, \href{https://mathoverflow.net/a/262422}{answer}, February 16, 2017;
J. Cigler, \href{https://mathoverflow.net/a/262469}{completion}, February 17, 2017;
T. Amdeberhan, \href{https://mathoverflow.net/a/262576}{independent proof}, February 18, 2017.
\bibitem{dlmf} NIST,
\emph{Digital Library of Mathematical Functions},
\href{https://dlmf.nist.gov/25.6.E2}{equation 25.6.2} and
\href{https://dlmf.nist.gov/4.19.E5}{equation 4.19.5},
accessed November 13, 2026.
\bibitem{bacher} R. Bacher,
\emph{Paperfolding and Catalan numbers},
\href{https://arxiv.org/abs/math/0406340v1}{arXiv:math/0406340v1},
June 17, 2004, Theorems 1.2(i) and 10.1(i).
\bibitem{guo} Y.-J. Guo,
\emph{On the regularity of the Hankel determinant sequence of the characteristic sequence of powers},
\href{https://arxiv.org/abs/1806.08729v1}{arXiv:1806.08729v1}, June 20, 2018,
Lemma 2(3)--(6), pp.~6--7; Advances in Applied Mathematics
\textbf{104} (2019), 100--116,
\href{https://doi.org/10.1016/j.aam.2018.12.001}{doi:10.1016/j.aam.2018.12.001}.
\bibitem{sun2019} Z.-W. Sun,
\emph{Open Conjectures on Congruences}, Nanjing University Journal of
Mathematics Biquarterly \textbf{36} (2019), no.~1, 1--99,
Conjecture 81(i) and Remark 81, p.~68,
\href{https://doi.org/10.3969/j.issn.0469-5097.2019.01.01}{doi:10.3969/j.issn.0469-5097.2019.01.01}.
\bibitem{edgar} G. A. Edgar,
\emph{The Ap\'ery Numbers as a Stieltjes Moment Sequence},
\href{https://arxiv.org/abs/2005.10733v2}{arXiv:2005.10733v2},
first released May 21, 2020, Theorem 1, p.~1.
\bibitem{adamczewski} T. Adamczewski,
\emph{OEIS Open: How many conjectures can language models turn into theorems?},
\href{https://arxiv.org/abs/2608.11941v2}{arXiv:2608.11941v2},
August 13, 2026, Table 1, p.~16, A228304.
\bibitem{community} D. Grinberg,
\href{https://math.stackexchange.com/a/2597700}{answer to \emph{A curious Hankel determinant}},
January 9, 2018, Theorem 1;
Random, \href{https://mathoverflow.net/a/513374}{MathOverflow answer},
July 19, 2026, the shifted power-of-two matrix.
\end{thebibliography}
\bigskip
\noindent\small\textcopyright\ 2026 Haobo Ma and Wenlin Zhang.
Licensed under Creative Commons Attribution 4.0 International (CC BY 4.0).
\end{document}
